
\paragraph{Flow-based generative models in physics applications.}

Flow-based generative models are increasingly used in physics for probabilistic forecasting, learning to sample the conditional distribution of the next state given the current one -- via noise added around neural-operator predictions \cite{lippe_pde-refiner_2023, gao_generative_2024}, or by direct conditional sampling with denoising diffusion \cite{kohl_benchmarking_2024}, flow matching \cite{fotiadis_stochastic_2024}, and stochastic interpolants \cite{chen_probabilistic_2024, mucke_physics-aware_2025}.

\paragraph{Posterior sampling with flow-based generative models.}
Generative models are increasingly used for inverse problems \cite{daras_survey_2024}, via training-free corrections of flow priors \cite{zhang_flow_2025, ahamed_dawn-si_2024, negrel_multitask_2025} or by conditioning a generator on measurements through the score decomposition $\nabla\log p_\tau(\bx\mid\obs)=\nabla\log p_\tau(\bx)+\nabla\log p_\tau(\obs\mid\bx)$ of \cite{song_score-based_2021}. For flow matching, OT-ODE \cite{pokle_training-free_2024} and FIG \cite{yan_fig_2024} guide the probability-flow ODE -- FIG with a measurement interpolant coinciding with ours \eqref{eq:interpolated_observation} -- while \cite{martin_pnp-flow_2025} folds the prior into a plug-and-play denoiser. The Gaussian likelihood surrogate these methods (and ours, Section~\ref{subsec:obs_interpolation}) rely on originates in diffusion guidance: DPS \cite{chung_diffusion_2023} uses the denoiser mean with uninflated covariance, $\Pi$GDM \cite{song_pseudoinverse_2023} inflates it. The likelihood-score coefficient in these ODE methods is fixed by the velocity--score conversion \cite{albergo_stochastic_2023, ma_sit_2024}, exactly the duality coefficient $\vscoef_\tau$ in our guided-ODE weight; our guided ODE generalises them to autoregressive DA, as the zero-diffusion member of the family containing the SDE samplers.

\paragraph{Data assimilation with generative models.}
Generative DA seeks to overcome the limitations of ensemble Kalman and particle filters. Early efforts include score-based DA \cite{rozet_score-based_2023}; for high-dimensional systems \cite{bao_ensemble_2024} introduced an ensemble score filter, extended to latent space for sparse observations by \cite{si_latent-ensf_2024}, with multimodal settings addressed by \cite{qu_deep_2024}. Most relevant is FlowDAS \cite{chen_flowdas_2025}, a stochastic-interpolant DA framework with a Monte-Carlo likelihood estimate; Appendix~\ref{appendix:comparison} details how our method differs from these and the methods above.
